\section{Datos de ejemplo utilizados}

Los datos son sintéticos y reproducibles. El generador usa la semilla 42 y fija
como referencia el 30 de junio de 2026 a las 15:00 UTC. El manifiesto guarda la
versión del generador, los conteos, el inventario de archivos, sus checksums y
los oráculos vectoriales.

\begin{table}[ht]
\centering
\small
\begin{tabular}{l r l}
\toprule
Grupo & Cantidad & Componentes principales \\
\midrule
Operación & 12 productos, 12 pedidos & 6 clientes, 28 líneas, 14 entregas y 4 incidencias \\
Documentación & 10 documentos & 14 versiones y 32 fragmentos \\
Vectores & 36 embeddings & 32 del modelo activo y 4 históricos \\
Acceso & 6 actores & 3 perfiles, 5 clases y 8 permisos \\
Interacción & 7 pares & 7 consultas, 7 respuestas y sus evidencias \\
Auditoría & 56 eventos & Cambios documentales, consultas, accesos y evidencia \\
\bottomrule
\end{tabular}
\caption{Escala principal del conjunto sintético.}
\end{table}

Los escenarios se definieron antes de variar valores con la semilla. Incluyen
un cliente sin pedidos, un pedido cancelado, entregas parciales, una demora con
incidencia, condiciones vencidas y futuras, todos los estados documentales, un
actor inactivo y accesos permitidos y denegados. Los 14 archivos Markdown son
ficticios; cada encabezado de segundo nivel produce un fragmento. Los oráculos
fijan como primeros resultados el fragmento 10 para la búsqueda vectorial y el
18 para la recuperación híbrida de Comercial/Compras.
